\begin{notechapter}
\chapter{Structured Prediction, Hidden Markov Models and the Viterbi Algorithm}\label{ch:hmm-viterbi}
\chaptersubtitle{Choosing the best of exponentially many tag sequences in linear time}

\begin{chapterintro}
Part-of-speech tagging assigns a tag (noun, verb, determiner, \ldots) to every word of a
sentence. The tags are not independent: a determiner is usually followed by a noun, and the
tag of ``duck'' depends on whether ``her'' is a pronoun or a possessive. Predicting the whole
sequence at once is an instance of \emph{structured prediction}, where the output is a
combinatorial object and the number of candidates grows exponentially with its size. For
chain-structured models, dynamic programming finds the best candidate exactly. This
chapter introduces structured prediction, the hidden Markov model (HMM) for tagging, and
the Viterbi and forward algorithms, with a proof, an example in which tagging word by word
fails, and verified code. You will learn
\begin{itemize}[nosep]
  \item what makes a prediction problem structured, and the two problems it raises: scoring
        and searching;
  \item the HMM, and its supervised estimation by counting;
  \item the Viterbi recursion, why it is exact, and its backtracking;
  \item the forward algorithm, which sums where Viterbi maximizes, and posterior marginals;
  \item why everything is done in log space.
\end{itemize}
\prerequisites{the bigram model and estimation by counting (\cref{ch:bigram-models}); Bayesian
networks (\cref{dkf:sec:dseparation}).}
\end{chapterintro}

\section{Structured prediction}\label{hmm-viterbi:sec:structured}

In classification the output is one label among $K$. In \emph{structured prediction} it is an object
with parts that constrain each other: a sequence of tags, a parse tree, a segmentation of an
image into regions, a translation. A tag sequence for a sentence of $T$ words, with $K$ possible tags,
has $K^T$ candidates: $45^{20}\approx10^{33}$ for a 20-word sentence and the 45 tags of the Penn Treebank. A
structured predictor needs
\begin{itemize}[nosep]
  \item a \emph{score} for every candidate $y$ given the input $x$, typically a probability $p(y\mid x)$ or $p(x, y)$, that
        decomposes over small parts of $y$ (pairs of neighbouring tags, for instance);
  \item a \emph{search} (inference, decoding) procedure that finds $\arg\max_yp(y\mid x)$ without enumerating the
        candidates, made possible by the decomposition;
  \item a \emph{learning} procedure, which usually needs sums over all candidates as well.
\end{itemize}
For chains, the decomposition is into neighbouring pairs, and search and summation are
done by dynamic programming. Hidden Markov models (this chapter) are the generative
model of this kind, conditional random fields (\cref{ch:crf}) the discriminative one. Other
structures use other algorithms (CKY for parse trees, graph cuts for some image models),
and when no exact algorithm is available, approximate search such as beam search, which
autoregressive sequence generators also use.

\section{Hidden Markov models}\label{hmm-viterbi:sec:hmm}

\begin{definition}[Hidden Markov model]\label{hmm-viterbi:def:hmm}
A hidden Markov model with states (tags) $\{1,\dots,K\}$ and observations (words) in a vocabulary defines the
joint distribution of a tag sequence $y_{1:T}$ and a word sequence $x_{1:T}$
\[
  p(x_{1:T}, y_{1:T}) = \pi_{y_1}B_{y_1}(x_1)\prod_{t=2}^TA_{y_{t-1}y_t}B_{y_t}(x_t),
\]
with initial probabilities $\pi_k$, transition probabilities $A_{kl} = p(y_t = l\mid y_{t-1} = k)$ and emission probabilities
$B_k(x) = p(x_t = x\mid y_t = k)$; each of $\pi$, each row of $A$ and each $B_k$ is a probability distribution.
\end{definition}

It is the Bayesian network of \cref{hmm-viterbi:fig:hmm}, the discrete-state counterpart of the
state-space model of \cref{ch:deep-kalman-filters}: the tags form a Markov chain (a bigram model of
tags), and each word depends only on its own tag. It is a \emph{generative} model: it describes how
the words are produced, and tagging uses Bayes' rule, $p(y\mid x) = p(x, y)/p(x)$.

\begin{figure}[ht]
\centering
\begin{tikzpicture}[
  v/.style={draw, circle, minimum size=8mm, inner sep=0pt, font=\small},
  ob/.style={v, fill=gray!30},
  ar/.style={-{Stealth[length=1.8mm]}, thick}]
  \foreach \t/\w in {1/I,2/saw,3/her,4/duck} {
    \node[v] (y\t) at (2.2*\t,0) {$y_{\t}$};
    \node[ob] (x\t) at (2.2*\t,-1.4) {$x_{\t}$};
    \node[font=\scriptsize] at (2.2*\t,-2.1) {``\w''};
    \draw[ar] (y\t) -- node[right, font=\scriptsize]{$B$} (x\t);
  }
  \foreach \t/\u in {1/2,2/3,3/4} \draw[ar] (y\t) -- node[above, font=\scriptsize]{$A$} (y\u);
  \node[font=\scriptsize] (pi) at (0.6,0) {$\pi$};
  \draw[ar] (pi) -- (y1);
\end{tikzpicture}
\caption{An HMM as a Bayesian network. The tags $y_t$ are hidden at test time, the words $x_t$ observed.}
\label{hmm-viterbi:fig:hmm}
\end{figure}

\paragraph{Supervised estimation.} With a tagged corpus, both tags and words are observed, the
log-likelihood is a sum of $\log\pi$, $\log A$ and $\log B$ terms, and by the argument of
\cref{bigram-models:prop:counts} its maximum is given by counts:
\[
  \hat A_{kl} = \frac{\#(k\to l)}{\#(k\to\cdot)},\qquad\hat B_k(x) = \frac{\#(k\text{ emits }x)}{\#(k)},\qquad\hat\pi_k = \frac{\#(\text{sentences starting with }k)}{\#(\text{sentences})},
\]
with smoothing for unseen events. The Baum--Welch algorithm, an instance of EM
(\cref{ch:em-algorithm}), is needed only when the tags are \emph{not} observed in the training data.

\section{The Viterbi algorithm}\label{hmm-viterbi:sec:viterbi}

Given words $x_{1:T}$, the most probable tag sequence maximizes $p(y_{1:T}\mid x_{1:T})$, equivalently $p(x_{1:T}, y_{1:T})$, since $p(x_{1:T})$ does
not depend on $y$. Define, for each time $t$ and tag $k$,
\[
  \delta_t(k) = \max_{y_{1:t-1}}p(x_{1:t},\ y_{1:t-1},\ y_t = k),
\]
the probability of the best tag sequence for the first $t$ words that ends with tag $k$ (jointly with
those words, not conditionally on them).

\begin{theorem}[Viterbi recursion]\label{hmm-viterbi:thm:viterbi}
$\delta_1(k) = \pi_kB_k(x_1)$, and for $t\geq2$
\[
  \delta_t(k) = \Bigl(\max_j\delta_{t-1}(j)\,A_{jk}\Bigr)B_k(x_t),\qquad\psi_t(k) = \arg\max_j\delta_{t-1}(j)\,A_{jk} .
\]
Then $\max_yp(x_{1:T}, y) = \max_k\delta_T(k)$, and a maximizing sequence is obtained by backtracking:
$y_T^* = \arg\max_k\delta_T(k)$ and $y_{t-1}^* = \psi_t(y_t^*)$ for $t = T,\dots,2$.
\end{theorem}

\begin{proof}
By the factorization, $p(x_{1:t}, y_{1:t}) = p(x_{1:t-1}, y_{1:t-1})\,A_{y_{t-1}y_t}B_{y_t}(x_t)$. Fix $y_t = k$ and split the maximum over $y_{1:t-1}$ into a
maximum over $y_{t-1} = j$ and over $y_{1:t-2}$. The factor $A_{jk}B_k(x_t)$ does not depend on $y_{1:t-2}$ and is nonnegative,
so it can be taken out of the inner maximum:
\[
  \delta_t(k) = \max_j\Bigl(\max_{y_{1:t-2}}p(x_{1:t-1}, y_{1:t-2}, y_{t-1} = j)\Bigr)A_{jk}B_k(x_t) = \max_j\delta_{t-1}(j)A_{jk}B_k(x_t) .
\]
For backtracking, let $y^*$ be the sequence it produces. By induction on $t$, the prefix $y_{1:t}^*$ attains $\delta_t(y_t^*)$:
it holds at $t = 1$, and $y_{t-1}^* = \psi_t(y_t^*)$ is the predecessor that attains the maximum in the recursion. At $t = T$,
$p(x, y^*) = \delta_T(y_T^*) = \max_k\delta_T(k) = \max_yp(x, y)$.
\end{proof}

This is the \emph{optimal substructure} of dynamic programming: the best path to $(t, k)$ extends one of the
best paths to $(t - 1, j)$, and only the $K$ values $\delta_{t-1}(\cdot)$ need to be kept, not the $K^{t-1}$ paths. The cost is
$O(TK^2)$ time and $O(TK)$ memory for the back-pointers, instead of $O(TK^T)$ for enumeration (Viterbi, \emph{Error
Bounds for Convolutional Codes and an Asymptotically Optimum Decoding Algorithm}, 1967).

\begin{example}[I saw her duck]\label{hmm-viterbi:ex:duck}
Tags P (pronoun), V (verb), D (determiner), N (noun). The initial probabilities are $\pi = (0.7, 0.05, 0.15, 0.1)$, the
transitions (row: from, column: to, in the order P, V, D, N)
\[
  A = \begin{bmatrix}0.05 & 0.8 & 0.05 & 0.1\\ 0.4 & 0.05 & 0.35 & 0.2\\ 0 & 0 & 0 & 1\\ 0.1 & 0.6 & 0.1 & 0.2\end{bmatrix},
\]
and the emissions of the four words (all others zero; each $B_k$ also gives probability to words
not in this sentence):
\begin{gather*}
  B_{\mathrm P}(\text{I}) = 0.4,\qquad B_{\mathrm V}(\text{saw}) = 0.1,\quad B_{\mathrm N}(\text{saw}) = 0.05,\\
  B_{\mathrm P}(\text{her}) = 0.3,\quad B_{\mathrm D}(\text{her}) = 0.25,\qquad B_{\mathrm V}(\text{duck}) = 0.03,\quad B_{\mathrm N}(\text{duck}) = 0.04 .
\end{gather*}
\emph{Recursion} (zeros omitted):
\begin{align*}
  \delta_1(\mathrm P) &= 0.7\cdot0.4 = 0.28,\\
  \delta_2(\mathrm V) &= 0.28\cdot0.8\cdot0.1 = 0.0224,\qquad \delta_2(\mathrm N) = 0.28\cdot0.1\cdot0.05 = 0.0014,\\
  \delta_3(\mathrm P) &= \max(0.0224\cdot0.4,\ 0.0014\cdot0.1)\cdot0.3 = 0.002688 &&(\psi_3(\mathrm P) = \mathrm V),\\
  \delta_3(\mathrm D) &= \max(0.0224\cdot0.35,\ 0.0014\cdot0.1)\cdot0.25 = 0.00196 &&(\psi_3(\mathrm D) = \mathrm V),\\
  \delta_4(\mathrm V) &= \max(0.002688\cdot0.8,\ 0.00196\cdot0)\cdot0.03 = 6.45\cdot10^{-5} &&(\psi_4(\mathrm V) = \mathrm P),\\
  \delta_4(\mathrm N) &= \max(0.002688\cdot0.1,\ 0.00196\cdot1)\cdot0.04 = 7.84\cdot10^{-5} &&(\psi_4(\mathrm N) = \mathrm D).
\end{align*}
\emph{Backtracking}: $y_4^* = \mathrm N$, $y_3^* = \psi_4(\mathrm N) = \mathrm D$, $y_2^* = \psi_3(\mathrm D) = \mathrm V$, $y_1^* = \mathrm P$: \emph{I/P saw/V her/D duck/N}, with
$p(x, y^*) = 7.84\times10^{-5}$.

\emph{Greedy tagging fails.} Choosing each tag in turn as the best continuation of the previous
choice picks P for ``her'' ($0.4\times0.3 = 0.12 > 0.35\times0.25 = 0.0875$) and then V for ``duck'', the reading \emph{I saw her
duck (dodge)}, with probability $6.45\times10^{-5}$. At $t = 3$ the pronoun is ahead, $\delta_3(\mathrm P) > \delta_3(\mathrm D)$, but the determiner
leads to a better completion. Viterbi keeps the best path to \emph{every} tag, so it does not have to
commit early.
\end{example}

\begin{figure}[ht]
\centering
\begin{tikzpicture}[
  c/.style={draw, rounded corners=2pt, minimum width=17mm, minimum height=7mm, font=\scriptsize, align=center},
  zc/.style={c, draw=gray!40, text=gray!60},
  ar/.style={-{Stealth[length=1.6mm]}, gray!60},
  best/.style={-{Stealth[length=2mm]}, very thick, bookpitfall}]
  \foreach \w [count=\t] in {I,saw,her,duck} \node[font=\small] at (2.7*\t,0.8) {``\w''};
  \foreach \k [count=\r] in {P,V,D,N} \node[font=\small] at (1.0,-1.0*\r) {\k};
  \node[c, fill=bookaccent!12] (P1) at (2.7,-1) {$0.28$};
  \node[zc] (V1) at (2.7,-2) {$0$}; \node[zc] (D1) at (2.7,-3) {$0$}; \node[zc] (N1) at (2.7,-4) {$0$};
  \node[zc] (P2) at (5.4,-1) {$0$}; \node[c, fill=bookaccent!12] (V2) at (5.4,-2) {$0.0224$};
  \node[zc] (D2) at (5.4,-3) {$0$}; \node[c, fill=bookaccent!12] (N2) at (5.4,-4) {$0.0014$};
  \node[c, fill=bookaccent!12] (P3) at (8.1,-1) {$0.002688$}; \node[zc] (V3) at (8.1,-2) {$0$};
  \node[c, fill=bookaccent!12] (D3) at (8.1,-3) {$0.00196$}; \node[zc] (N3) at (8.1,-4) {$0$};
  \node[zc] (P4) at (10.8,-1) {$0$}; \node[c, fill=bookaccent!12] (V4) at (10.8,-2) {$6.45\cdot10^{-5}$};
  \node[zc] (D4) at (10.8,-3) {$0$}; \node[c, fill=bookexample!25] (N4) at (10.8,-4) {$7.84\cdot10^{-5}$};
  \draw[ar] (P1) -- (N2); \draw[ar] (V2) -- (P3); \draw[ar] (P3) -- (V4);
  \draw[best] (P1) -- (V2); \draw[best] (V2) -- (D3); \draw[best] (D3) -- (N4);
\end{tikzpicture}
\caption{The trellis of \cref{hmm-viterbi:ex:duck}: $\delta_t(k)$ in each cell, the back-pointers $\psi$ of the non-zero cells as
arrows. The best path (red) goes through $\delta_3(\mathrm D)$, which is smaller than $\delta_3(\mathrm P)$ at time $3$.}
\label{hmm-viterbi:fig:trellis}
\end{figure}

\section{The forward algorithm and posterior marginals}\label{hmm-viterbi:sec:forward}

Replacing the maximum by a sum gives the probability of the words, summed over all tag
sequences. With $\alpha_t(k) = p(x_{1:t}, y_t = k)$,
\[
  \alpha_1(k) = \pi_kB_k(x_1),\qquad\alpha_t(k) = \Bigl(\sum_j\alpha_{t-1}(j)A_{jk}\Bigr)B_k(x_t),\qquad p(x_{1:T}) = \sum_k\alpha_T(k),
\]
by the same argument as \cref{hmm-viterbi:thm:viterbi}, with the distributivity of the product over
the sum in place of its distributivity over the maximum. A backward recursion,
$\beta_T(k) = 1$, $\beta_t(k) = \sum_lA_{kl}B_l(x_{t+1})\beta_{t+1}(l) = p(x_{t+1:T}\mid y_t = k)$, then gives the \emph{posterior marginals}
\[
  p(y_t = k\mid x_{1:T}) = \frac{\alpha_t(k)\,\beta_t(k)}{p(x_{1:T})} .
\]
In the example, $p(x_{1:4}) = 1.5624\times10^{-4}$, so the Viterbi sequence has posterior probability $7.84/15.624 = 0.502$:
half of the probability mass is spread over other taggings. The posterior marginal of D at $t = 3$
is $0.511$ and that of N at $t = 4$ is $0.581$. Choosing the most probable tag at each position separately
(\emph{posterior decoding}) maximizes the expected number of correct tags; Viterbi maximizes the
probability that the whole sequence is correct. The two can disagree, and posterior decoding
can even produce a sequence of probability $0$ (\cref{hmm-viterbi:exr:4}). The forward--backward
recursions are also what the E-step of Baum--Welch and the training of CRFs
(\cref{ch:crf}) need.

\section{Log space and implementation}\label{hmm-viterbi:sec:code}

The probabilities $\delta_t(k)$ shrink geometrically with $t$ (here to $10^{-5}$ after four words) and underflow
in floating point for sentences of a few hundred words. Since the logarithm is increasing, the
Viterbi recursion can be run on $\log\delta$, with products replaced by sums and the same maximizers:
\[
  \log\delta_t(k) = \max_j\bigl(\log\delta_{t-1}(j) + \log A_{jk}\bigr) + \log B_k(x_t) .
\]
For the forward algorithm, the logarithm of a sum is computed with the log-sum-exp function,
$\log\sum_je^{a_j} = m + \log\sum_je^{a_j - m}$, $m = \max_ja_j$.
\begin{lstlisting}[language=Python]
import numpy as np
from scipy.special import logsumexp   # handles -inf entries correctly

def viterbi(log_pi, log_A, log_B):        # log_B[t, k] = log B_k(x_t)
    T, K = log_B.shape
    delta = log_pi + log_B[0]
    psi = np.zeros((T, K), dtype=int)
    for t in range(1, T):
        scores = delta[:, None] + log_A      # scores[j, k] = delta_{t-1}(j) + log A_jk
        psi[t] = scores.argmax(axis=0)
        delta = scores.max(axis=0) + log_B[t]
    path = [int(delta.argmax())]
    for t in range(T - 1, 0, -1):            # backtracking
        path.append(int(psi[t, path[-1]]))
    return path[::-1], delta.max()

def log_forward(log_pi, log_A, log_B):
    a = log_pi + log_B[0]                    # log alpha_1
    for t in range(1, log_B.shape[0]):
        a = logsumexp(a[:, None] + log_A, axis=0) + log_B[t]
    return logsumexp(a)                      # log p(x_1:T)
\end{lstlisting}
On the example, with \texttt{np.log} turning the zero probabilities into $-\infty$, the functions return
the path P, V, D, N and $\log p(x) = \log(1.5624\times10^{-4})$, which a brute-force enumeration of the $4^4 = 256$ tag
sequences confirms. A hand-written log-sum-exp must treat the case where all the terms are
$-\infty$ (the shift $m$ is then $-\infty$ and $a_j - m$ is undefined); \texttt{scipy.special.logsumexp} does.

\begin{chaptersummary}
\begin{itemize}
  \item Structured prediction chooses among exponentially many outputs; it needs a score that
        decomposes over small parts and an algorithm that exploits the decomposition.
  \item An HMM is a Markov chain of tags emitting one word each; with tagged data, its parameters
        are conditional frequencies (Baum--Welch is for untagged data).
  \item Viterbi: $\delta_t(k) = \max_j\delta_{t-1}(j)A_{jk}B_k(x_t)$ with back-pointers gives the most probable tag sequence exactly, in
        $O(TK^2)$; greedy tagging can fail because it commits early.
  \item The forward algorithm replaces the max by a sum and gives $p(x)$; with the backward recursion it
        gives posterior marginals.
  \item Implementations work with logarithms (and log-sum-exp) to avoid underflow.
\end{itemize}
\end{chaptersummary}

\section*{Exercises}
\addcontentsline{toc}{section}{Exercises}

\begin{exercise}\label{hmm-viterbi:exr:1}
Count the multiplications of Viterbi and of brute-force enumeration for $T = 20$ and $K = 45$.
\solutionlink{hmm-viterbi:exr:1}
\end{exercise}

\begin{exercise}\label{hmm-viterbi:exr:2}
In \cref{hmm-viterbi:ex:duck}, compute $\alpha_3$, $\alpha_4$ and $p(x_{1:4})$, and check the posterior probability $0.502$ of the Viterbi
sequence.
\solutionlink{hmm-viterbi:exr:2}
\end{exercise}

\begin{exercise}\label{hmm-viterbi:exr:3}
Estimate $\pi$, $A$ and $B_{\mathrm N}$ by counting from the tagged sentences \emph{the/D dog/N runs/V}, \emph{the/D dogs/N
run/V}, \emph{dogs/N run/V}.
\solutionlink{hmm-viterbi:exr:3}
\end{exercise}

\begin{exercise}\label{hmm-viterbi:exr:4}
Show that with two tags and two time steps, if each posterior marginal has a unique maximum,
the tags chosen separately by posterior decoding form a sequence of positive probability. Then construct an HMM with three tags and two
time steps in which they form a sequence of probability $0$.
\hint{Use a transition probability equal to $0$, and uninformative emissions.}
\solutionlink{hmm-viterbi:exr:4}
\end{exercise}

\begin{exercise}\label{hmm-viterbi:exr:5}
Show that the Viterbi recursion remains correct when $\max$ and products are replaced by $\max$ and
sums of logarithms, and explain why it would not remain correct if the logarithm were replaced
by a decreasing function.
\solutionlink{hmm-viterbi:exr:5}
\end{exercise}
\end{notechapter}
